\documentclass[11pt]{article}
\usepackage[utf8]{inputenc}
\usepackage{amsmath,amssymb}
\usepackage{graphicx}
\usepackage{booktabs}
\usepackage{hyperref}
\usepackage[margin=1in]{geometry}
\usepackage{natbib}
\bibliographystyle{unsrtnat}

\title{Stable Population Structure in Europe Since the Iron Age Despite High Individual Mobility}
\author{}
\date{}

\begin{document}
\maketitle

\begin{abstract}
Ancient DNA studies have revealed dramatic population transformations during European prehistory, yet the genetic dynamics of the historical period (the last 3,000 years) remain poorly characterized. Here we report 204 new whole genomes from 53 archaeological sites across 18 countries spanning the Iron Age through the Early Modern period, including the first historical-period genomes from Armenia, Algeria, Austria, and France. Despite finding that at least 7\% of historical individuals carry ancestry uncommon in their sampling region---indicating extensive cross-Mediterranean and cross-continental contacts---we demonstrate that overall population genetic structure has remained remarkably stable and geographically concordant from the Iron Age to the present. Stepping-stone migration simulations reveal that the observed level of individual dispersal under standard panmictic models would collapse population structure within a few generations, implying that effective migration rates were substantially lower than apparent dispersal. We propose a transient dispersal hypothesis in which improved transportation networks and Roman Empire mobilization moved people across vast distances without proportional gene flow, reconciling high individual mobility with persistent population structure.
\end{abstract}

\section{Introduction}

The application of ancient DNA (aDNA) to European prehistory has uncovered successive waves of population change: the Neolithic farmer expansion beginning around 7,500~BCE \citep{haak2015massive,lazaridis2014ancient}, followed by large-scale migrations of Steppe pastoralists during the Bronze Age \citep{allentoft2015population,haak2015massive}. By the end of the Bronze Age, the three major ancestry components of present-day Europeans---Western Hunter-Gatherer (WHG), Neolithic Farmer, and Bronze Age Steppe Herder---were largely established \citep{lazaridis2014ancient}. However, remarkably little is known about genetic dynamics during the subsequent historical period spanning the last 3,000 years, despite this era encompassing the rise and fall of empires, mass migrations, trade networks, and colonial expansion \citep{antonio2019ancient,olalde2019genomic}.

Individual regional studies have revealed high genetic heterogeneity in specific historical contexts. Ancient genomes from Imperial Rome showed extraordinary diversity, with individuals carrying ancestries from across the Mediterranean and beyond \citep{antonio2019ancient}. Similarly, studies of the Iberian Peninsula documented complex admixture patterns during historical periods \citep{olalde2019genomic}. Yet no comprehensive multi-regional assessment has examined whether these patterns of individual mobility translate into detectable shifts in overall population structure.

This gap raises a fundamental paradox. Standard population genetics theory predicts that even modest levels of migration between populations will rapidly homogenize allele frequencies \citep{wright1931evolution,slatkin1987gene}. If 7\% or more of individuals in any generation are migrants from distant regions, population structure should erode within tens of generations. Yet present-day European populations exhibit clear geographic structure that correlates strongly with geography \citep{novembre2008genes,lao2008correlation}. Either historical mobility was lower than individual case studies suggest, or some mechanism decoupled individual movement from effective gene flow.

Here we address these questions through the largest multi-regional aDNA study of the European historical period to date. We generated 204 new whole genomes from archaeological sites spanning 18 countries and the full historical period, analyzed alongside 2,033 present-day genomes, 1,998 prehistoric genomes, and 764 published historical genomes. Our results demonstrate that population structure has been remarkably stable since the Iron Age despite substantial individual-level dispersal, and we propose a transient dispersal model to explain this apparent paradox.

\section{Results}

\subsection{A multi-regional ancient genome dataset spanning the historical period}

We generated whole-genome shotgun data from 204 individuals excavated from 53 archaeological sites in 18 countries (Figure~1A--B). DNA was extracted primarily from petrous bones ($n = 203$) and teeth ($n = 1$), with libraries treated using partial uracil-DNA glycosylase (UDG) to reduce ancient DNA damage artifacts while preserving informative damage patterns at terminal positions. Sequencing achieved a median depth of 0.92$\times$ (range 0.16--2.38$\times$), and pseudohaploid genotypes were called on the 1240k SNP panel, yielding a median of 685,058 SNPs per sample (range 167,000--1,029,345). Radiocarbon dating was performed on 126 samples to establish precise temporal placement.

The dataset was organized into 14 geographic categories across seven time periods: Mesolithic and Neolithic (10000--3500~BCE), Copper Age (3500--2300~BCE), Bronze Age (2300--1000~BCE), Iron Age (1000~BCE--1~CE), Imperial Rome and Late Antiquity (1--700~CE), Medieval Ages and Early Modern (700--1900~CE), and present-day (1900~CE onward). Critically, this study reports the first historical-period genomes from Armenia, Algeria, Austria, and France, substantially expanding geographic coverage of the historical aDNA record.

\subsection{Population structure mirrors geography across all historical periods}

Principal component analysis (PCA) projection of ancient genomes onto present-day European genetic variation revealed a striking finding: population structure closely mirrored geography across all historical time periods (Figure~2, Supplementary Figure~2). Using smartpca (v1600) with 829 present-day individuals and 480,712 SNPs, we projected ancient individuals into a shared genetic space that reflects the well-established correlation between genetic and geographic distances in Europe \citep{novembre2008genes}.

During the Iron Age, individuals from each region clustered near their corresponding present-day populations. This geographic concordance persisted through the Imperial Roman and Late Antique periods, the Medieval Ages, and into the present day. Despite the massive political, cultural, and demographic upheavals of these two millennia, the overall genetic landscape of Europe remained remarkably stable.

Regional case studies reinforced this finding. In Armenia, we identified two genetically homogeneous clusters that were temporally distinct, each projecting to positions consistent with present-day Caucasus populations (Figure~2). Southeastern European samples showed high heterogeneity during the Imperial Roman period, with multiple distinct genetic clusters coexisting in the same time and place (Figure~3). Western European samples similarly revealed heterogeneity with evidence for cross-Mediterranean contacts, particularly during the Imperial period (Figure~4). In all regions, however, the central tendency of each population remained geographically anchored.

\subsection{At least 7\% of historical individuals are ancestry outliers}

To quantify individual-level mobility, we developed an outlier detection pipeline combining pairwise qpAdm modeling with clustering. For each individual, we assessed whether their ancestry was consistent with their sampling region or instead matched a distant population. An individual was classified as an outlier if their ancestry was better modeled as originating from a non-local source population.

Across all regions and time periods combined, at least 7\% of historical individuals were classified as ancestry outliers (Figure~5A). Italy during the Imperial Roman period showed the highest outlier proportions, consistent with Rome's role as a cosmopolitan hub drawing people from across the Mediterranean and beyond \citep{antonio2019ancient}. Southeastern Europe similarly showed high heterogeneity, while Western Europe displayed moderate levels and Armenia showed lower heterogeneity.

One-component qpAdm source tracing of outlier clusters revealed cross-Mediterranean and cross-continental distances between individuals and their putative source populations (Figure~5B--C). No significant sex bias was detected among outliers ($p = 0.4117$ for outlier versus non-outlier male/female proportions; $p = 0.633$ for outliers with and without identified sources), suggesting that both men and women participated in long-distance mobility during the historical period (Supplementary Figure~5-1).

Quality control analyses confirmed that outlier classification was not driven by differential SNP coverage. Median SNP counts showed no significant correlation with cluster size, and outlier clusters did not differ significantly in SNP coverage from non-outlier clusters (Supplementary Figure~2-3).

\subsection{Simulations reveal a mobility-stability paradox}

The coexistence of 7\%+ individual dispersal with stable population structure presents a paradox under standard population genetics models. To test this formally, we implemented stepping-stone migration simulations with local panmixia, parameterized by the observed outlier proportions.

The simulations demonstrated that under standard assumptions---where migrants reproduce at rates equivalent to the local population---the observed level of cross-regional dispersal would rapidly collapse population structure, homogenizing allele frequencies across Europe within tens of generations (Figure~5, Supplement~3). Since population structure has demonstrably persisted over 100+ generations from the Iron Age to the present, effective migration rates must be substantially lower than apparent dispersal rates.

\subsection{A transient dispersal hypothesis}

We propose that the discrepancy between high individual mobility and low effective gene flow reflects transient dispersal: individuals traveled long distances---facilitated by Roman roads, maritime routes, military mobilization, and trade networks---but did not proportionally contribute to the gene pool of their destination populations. Several mechanisms could produce this effect. Military deployments and commercial ventures may have been predominantly male and temporary, with individuals returning to their regions of origin. Enslaved people and laborers moved to distant locations may have had reduced reproductive success. Urban migrants may have experienced higher mortality rates than the general population \citep{scheidel2001death}.

Analysis of travel routes and times across the Roman Empire using the ORBIS model \citep{scheidel2014orbis} illustrates that the fastest summer routes for civilians---combining road, river, coastal, and open sea travel---could traverse the Mediterranean in weeks to months (Supplementary Figure~5-3). This represents a dramatic compression of travel timescales compared to prehistoric migrations, which unfolded over hundreds of years. The Roman transportation network thus enabled a novel form of dispersal that was rapid, extensive, and---critically---decoupled from gene flow.

\section{Discussion}

Our analysis of 204 new historical-period genomes alongside thousands of published ancient and modern genomes reveals a fundamental feature of European population history: the remarkable stability of population genetic structure over the last 3,000 years. This stability persists despite unambiguous evidence for extensive individual mobility, including at least 7\% ancestry outliers indicating cross-regional and cross-Mediterranean contacts.

This finding has important implications for how we interpret ancient DNA evidence of individual mobility. The presence of a non-local individual at an archaeological site---however distant their ancestry---does not necessarily imply lasting gene flow. Our simulations show that the observed mobility levels, if translated directly into effective migration, would have erased the very structure we observe today. The transient dispersal model provides a framework for understanding how empires, trade networks, and improved transportation can move people without moving populations.

The three-component ancestry model established during the Bronze Age---Western Hunter-Gatherer, Neolithic Farmer, and Bronze Age Steppe Herder---appears to have been essentially locked in by the Iron Age, with subsequent historical events producing only minor regional adjustments rather than continental-scale restructuring. This contrasts sharply with the dramatic turnovers documented during the Neolithic and Bronze Age transitions \citep{haak2015massive,allentoft2015population}.

Several limitations should be noted. Low-coverage shotgun sequencing (median 0.92$\times$) limits the detection of fine-scale population structure, and the 1240k SNP panel, while widely used, captures only a fraction of genomic variation. The outlier detection threshold is necessarily arbitrary, and the true proportion of non-local individuals may differ from our 7\% estimate. Future work with higher-coverage genomes and whole-genome analyses may reveal subtler patterns of admixture and gene flow.

\section{Methods}

\subsection{Sample collection and DNA extraction}
DNA was extracted from petrous bones ($n = 203$) and teeth ($n = 1$) following established protocols for ancient DNA recovery. Libraries were prepared with partial UDG treatment to reduce cytosine deamination artifacts while retaining informative damage patterns at read termini.

\subsection{Sequencing and genotyping}
Whole-genome shotgun sequencing was performed, achieving a median depth of 0.92$\times$ (range 0.16--2.38$\times$). Pseudohaploid genotypes were called on the 1240k SNP panel. Radiocarbon dating was performed on 126 samples.

\subsection{Population structure analysis}
PCA was performed using smartpca (v1600) with 829 present-day individuals and 480,712 SNPs, with 5 outlier iterations using 10 PCs for outlier removal. Ancient genomes were projected onto the present-day PCA space using least-squares projection (\texttt{lsqproject = YES}).

\subsection{Ancestry modeling}
qpAdm was used for ancestry decomposition. Outlier individuals were identified through individual pairwise qpAdm modeling followed by clustering. Source populations for outlier clusters were traced using one-component qpAdm models.

\subsection{Migration simulations}
Stepping-stone models with local panmixia were implemented to test whether observed dispersal rates could maintain population structure under standard population genetics assumptions. Migration rates were parameterized from observed outlier proportions.

\subsection{Data availability}
New sequences are deposited in the European Nucleotide Archive under accessions PRJEB53565 (raw) and PRJEB53564 (mapped). Previously reported sequences are available under PRJEB49419. Published comparative data were obtained from the Allen Ancient Data Resource (AADR).

\bibliography{references}

\end{document}
